\paragraph{Setup.}
We begin with a classic LRH setting: words are defined by features
$z=(z_1,\ldots,z_m)$; for example, \emph{royal} and \emph{female}
are features active for the word \emph{queen}. We allow signed
features, $z_i\in\{-1,0,1\}$, with zero denoting absence. The network
receives a word $z$ as input and predicts an associated signed
feature value, $\hat p(i^\pm\mid z)$.

The input distribution $p(z)$ determines which features co-occur
and how often they are encountered. A state with $k=\|z\|_0$ active
features is sparse when $k$ is small relative to $m$. For singleton
states, we write
\begin{equation}
p(z=e_i)=p(z=-e_i)=\frac{\pi_i}{2},
\qquad
\sum_{i=1}^m\pi_i=1.
\label{eq:frequency}
\end{equation}
We distinguish input frequency $\pi_i$ from predictive strength
$\alpha_i$: frequency determines how often a prediction contributes
to the loss, while strength determines how much an active feature
changes the conditional output distribution.

\subsection{Predicting resource-constrained representations.}

With enough dimensions, a representation can reproduce the target
logits. With fewer dimensions, it must discard directions or let
features overlap. Related work studies the capacity of linear
feature representations~\cite{garg} and their spectral
geometry~\cite{spectral}. We study this tradeoff behind
superposition~\cite{elhage} under the same softmax readout, first
for one-dimensional features and then for conceptual spaces.

Feature $i$ has outputs $i^+$ and $i^-$, with target and model logits
\begin{equation}
\ell^\star(i^\pm\mid z)=\pm\alpha_i z_i,
\qquad
\hat\ell(i^\pm\mid z)=E(z)^\top D_{i^\pm},
\qquad
E(z),D_{i^\pm}\in\R^d.
\label{eq:model}
\end{equation}
One softmax over all $2m$ outputs gives the predictions. The encoder
and decoder are untied; in particular, we do not require
$D_{i^-}=-D_{i^+}$. For $d<m$, we fit
\begin{equation}
\min_{E,D}\ \E_z\,
\CE[p_\star(\cdot\mid z),\hat p(\cdot\mid z)].
\label{eq:objective}
\end{equation}
This is exponential-family PCA: low-rank logits are fitted by
predictive loss rather than squared error~\cite{collins}.
Truncating the target SVD need not solve it.

We first consider weak and sharp predictions, then ask how
unequal input frequencies affect the finite-strength solution.
For the uniform sparsity comparison, we choose $k$ active
coordinates uniformly, with independent fair signs.
\begin{center}
\renewcommand{\arraystretch}{1.12}
\begin{tabular}{@{}lcc@{}}
\toprule
 & Weak logits & Strong logits\\
\midrule
Co-active & Reconstruct the logits & Preserve several outputs\\
Singleton & Reconstruct the logits & Distinguish one alternative\\
\bottomrule
\end{tabular}
\end{center}

\subsection{Weak logits: keep the strongest directions.}

Take the dense case, $k=m$, and $\alpha_i=\varepsilon a_i$,
with fixed $a_i>0$ and $\varepsilon\to0$. Predictions are near
uniform, and cross-entropy becomes squared logit error to leading
order. Write $L_0^\star=U\Sigma V^\top$ for the target-logit table
divided by $\varepsilon$, with rows weighted by $\sqrt{P(z)}$.
The limiting problem is
\begin{equation}
\boxed{
\min_{\rank(\hat L_0)\le d}
\|L_0^\star-\hat L_0\|_F^2
\quad\Longrightarrow\quad
\hat L_0=U_d\Sigma_dV_d^\top.
}
\label{eq:svd}
\end{equation}
The independent feature axes are already singular directions,
with singular values $\sqrt2\,a_i$. If $a_1>\cdots>a_m$, the model
keeps the strongest $d$ features and drops the rest:
\begin{equation}
\hat\ell(i^\pm\mid z)
\approx
\pm\alpha_i z_i\,\mathbf1\{i\le d\}.
\label{eq:truncation}
\end{equation}

Does sparsity change this geometry? Under uniform $k$-active
sampling, reducing $k$ rescales all singular values by the same
factor $\sqrt{k/m}$, leaving the leading spectral allocation
unchanged. Unequal singleton frequencies instead give singular
values $\sqrt{2\pi_i}\,a_i$, so modes are ordered by
$\pi_i a_i^2$. With a common strength $\alpha_i=\alpha$, the
weak solution therefore retains the most frequent features.
Equal mode strengths leave the retained subspace degenerate.
Appendix~\ref{app:weak} derives these results.

\subsection{Sharp predictions: preserve distinctions.}

\subsubsection{Singleton predictions.}

Now take $z\sim\Unif\{\pm e_i\}_{i=1}^m$ and equal strengths
$\alpha\to\infty$, with $m$ fixed. One output dominates. Losing
logits need not be reconstructed accurately if their probabilities
remain negligible. In particular, the opposite value of the active
feature no longer needs to be distinguished from other incorrect
outputs at leading order.

This connects the task to the categorical margin geometry studied
in neural collapse~\cite{lu,jiang}. With encoder and decoder
directions normalized, the large-scale one-hot objective gives
the Softmax-code problem~\cite{jiang}:
\begin{equation}
\boxed{
\max_{\{D_y:\,\|D_y\|=1\}}
\ \min_y\
\operatorname{dist}\!\left(
D_y,\operatorname{conv}\{D_{y'}:y'\ne y\}
\right).
}
\label{eq:packing}
\end{equation}
Here $y$ ranges over the $2m$ signed output labels. Each output
must be separated from the convex hull of its competitors.
In two dimensions, the solution is a regular polygon:
\begin{equation}
\theta_a=\frac{2\pi(a-1)}{2m},
\qquad
D_{y_a}=(\cos\theta_a,\sin\theta_a)^\top,
\qquad a=1,\ldots,2m.
\label{eq:circle}
\end{equation}
The signed labels may be assigned to these vertices without
requiring the two values of a feature to be antipodal.

This polygon also supplies a feasible solution to our
rank-constrained predictive problem: scaling the corresponding
vertex encodings proportionally to $\alpha$ drives predictive KL
to zero (Appendix~\ref{app:code}). All $m$ features remain
distinguishable in two dimensions, although the target logits
have rank $m$. The regular spacing is the normalized max-margin
solution; rank alone does not uniquely determine the angles.

\subsubsection{Several active features.}

When $k>1$ features are active, several outputs must survive
together. Let
\[
S(z)=\{i^+:z_i=1\}\cup\{i^-:z_i=-1\}.
\]
For equal strengths and $|S(z)|=k$,
\begin{equation}
p_\star(y\mid z)
\longrightarrow
\frac1k\,\mathbf1\{y\in S(z)\}.
\label{eq:sharp-support}
\end{equation}
A fixed decoder can represent this limit exactly when there is
a direction $n_S$ satisfying
\begin{equation}
n_S^\top D_y=c_S \quad (y\in S),
\qquad
n_S^\top D_y<c_S \quad (y\notin S).
\label{eq:face}
\end{equation}
Geometrically, the labels in $S$ must be exactly those lying on
an exposed face of the decoder's convex hull~\cite{rauh}.
Scaling $E(z)=\lambda n_{S(z)}$ then suppresses all other outputs.
Singleton prediction requires selectable vertices; joint
prediction requires selectable faces.

There is a simple construction supporting every small
conjunction. Assign distinct numbers $t_y$ to the $2m$ output
labels and place their decoder vectors on the moment curve:
\begin{equation}
D_y=(t_y,t_y^2,\ldots,t_y^{2k})^\top.
\label{eq:moment-decoder}
\end{equation}
Every subset of at most $k$ points forms a face of this cyclic
polytope, the property called
$k$-neighborliness~\cite{rauh,donoho-tanner}. To see the
corresponding encoder, form
\[
q_S(t)=\prod_{y\in S}(t-t_y)^2.
\]
Its negative nonconstant coefficients give an encoding whose
logit at $y$ is $c_0-q_S(t_y)$: all selected outputs tie, and
all others lie below them.

Thus, for fixed finite $m$, every nonempty signed support of size
at most $k$ admits a sharp representation in $2k$ dimensions.
Scaling these encodings proportionally to $\alpha$ makes
predictive KL vanish; the empty state is represented by $E(0)=0$.
The encoder can depend nonlinearly on the entire active set.
An additive predictive task therefore need not have an additive
compressed representation. Appendix~\ref{app:multi} gives the
construction and proof.

\subsection{Unequal frequencies: when does a dropped feature enter?}

Return to singleton inputs, with common finite strength
$\alpha_i=\alpha$ and frequencies
$\pi_1>\cdots>\pi_m$. The weak solution retains features
$1,\ldots,d$ and drops the rest. At finite $\alpha$, this
coordinate-truncated solution remains a stationary model:
retained features are predicted exactly, and omitted inputs
receive uniform predictions. We ask when it becomes advantageous
for an omitted feature to share a retained direction.

Take $i\le d$ and $j>d$. Allow their encoder and decoder
coefficients to change independently, so that the positive-state,
positive-output logit block becomes
\begin{equation}
\alpha
\begin{pmatrix}
1 & v\\
u & uv
\end{pmatrix}.
\label{eq:mixing-perturbation}
\end{equation}
The negative-state and negative-output entries follow by sign
reversal for this perturbation. The block remains rank one.
The off-diagonal terms introduce interference, while the
$uv$ term begins to encode feature $j$.

The second variation has a negative direction exactly when
\begin{equation}
\boxed{
\frac{\pi_j}{\pi_i}>\rho_c(\alpha,m),
\qquad
\rho_c(\alpha,m)=
\frac{\alpha^2(m-1+\cosh\alpha)}
{m\sinh^2\alpha}.
}
\label{eq:frequency-threshold}
\end{equation}
This is an instability within an admissible untied perturbation,
so whenever the inequality holds, the coordinate-truncated
solution cannot be a local minimum of the full model.
Appendix~\ref{app:mixed} derives the second variation.

For small $\alpha$,
\begin{equation}
\rho_c(\alpha,m)
=
1-C_m\alpha^2+O(\alpha^4),
\qquad
C_m=\frac13-\frac1{2m}.
\label{eq:threshold-expansion}
\end{equation}
Only features sufficiently close in frequency can initially
share a direction. The first retained--omitted pair to become
unstable is $(d,d+1)$.

For Zipf frequencies $\pi_i\propto i^{-\beta}$, $\beta>0$, this
first instability occurs when
\begin{equation}
\left(\frac{d}{d+1}\right)^\beta
>
\rho_c(\alpha,m).
\label{eq:zipf-onset}
\end{equation}
More generally, a retained feature $i$ can participate in an
unstable pair only if
$i>(d+1)\rho_c^{1/\beta}$, and an omitted feature $j$ can
participate only if $j<d\rho_c^{-1/\beta}$. These conditions
identify a band around the spectral cutoff. For large $d$,
its boundaries are approximately
\begin{equation}
h_{\rm lin}\approx d\,\rho_c^{1/\beta},
\qquad
t_{\rm lin}\approx d\,\rho_c^{-1/\beta},
\label{eq:zipf-band}
\end{equation}
up to integer endpoints. Where a continuum approximation is
appropriate, its small-$\alpha$ width is
\[
t_{\rm lin}-h_{\rm lin}
\approx
\frac{2C_m}{\beta}\,d\alpha^2.
\]
These are boundaries of the linear instability, not a general
formula for the final nonlinear optimum. They suggest a
dedicated head, a superposed body near the capacity boundary,
and an omitted tail.

A small linear model tests this prediction directly. Take
$m=12$, $d=4$, $\pi_i\propto i^{-1}$, and $\alpha=1.15$.
Then $\rho_c\approx0.6954$, while
\[
\frac{\pi_5}{\pi_4}=0.8>\rho_c,
\qquad
\frac{\pi_6}{\pi_4}=\frac23<\rho_c,
\qquad
\frac{\pi_5}{\pi_3}=0.6<\rho_c.
\]
Only the retained--omitted pair $(4,5)$ is unstable by this
criterion. Optimizing the unrestricted rank-four logit
factorization from random initialization recovers three
dedicated feature directions, features $4$ and $5$ sharing the
fourth, and negligible representation of features $6,\ldots,12$.
The observed head/body/tail split therefore agrees with the
local prediction, without prescribing those groups in the model.

\subsection{What changes when features are correlated?}

Correlation changes which predictive combinations are worth
retaining. For mean-zero $z$ with covariance $\Sigma_z$, weak
prediction keeps the leading left singular vectors of
$\diag(a)\Sigma_z^{1/2}$ rather than necessarily individual
feature axes. The preceding weak-regime selection therefore
applies to predictive modes; their geometry can then be
expressed in the original feature coordinates.

A separate mixed-coordinate example makes this transformation
explicit. Let $z=Az'$, with $A$ invertible and $z'$ uniform on
signed singletons; outputs label the patterns $z'$. Write the
linear latent encoding as $Hz'$ and the observed-coordinate
encoding as $E(z)=Bz$. The encoder undoes the mixing:
$B=HA^{-1}$. Consequently,
\begin{equation}
\boxed{
B^\top B=A^{-\top}H^\top H A^{-1}.
}
\label{eq:correlated-gram}
\end{equation}
This identity uses the embedding geometry, without tying
embeddings to decoder vectors.

In the uncompressed isotropic reference $H^\top H=cI_m$,
we have $\Sigma_z=AA^\top/m$, giving
\begin{equation}
\boxed{
B^\top B=\frac{c}{m}\Sigma_z^{-1},
\qquad
\cos(B_i,B_j)=-\rho_{ij\,\cdot\,\mathrm{rest}}.
}
\label{eq:partial}
\end{equation}
Here $\rho_{ij\,\cdot\,\mathrm{rest}}$ is correlation after
linearly controlling for the other coordinates. Positive
partial correlation gives opposing encoder coefficients:
separating contributions requires canceling shared variation.
These are jointly fitted coefficients, not marginal activation
contrasts. Under compression, $H^\top H$ contains the
capacity-selected mode geometry instead of an isotropic metric.
Equation~\eqref{eq:correlated-gram} still applies, but the simple
partial-correlation identity need not hold
(Appendix~\ref{app:corr}).

\subsection{When a feature occupies a space}

A feature need not occupy one direction: an angle requires a plane. A rank constraint can remove some of its predictive modes or force its space to overlap others. In the weak regime, SVD may discard a feature's weak axis while retaining its strong one.

To study overlap among retained modes, collect feature $i$'s decoder directions into a block $D_i$, with internal coordinates $z_i$. Under $E(z)=\sum_jD_jz_j$, block $j$ contributes $D_i^\top D_jz_j$ to block $i$'s readout. Thus $D_i^\top D_j$ measures interference between spaces.

Consider two planes in three dimensions. Activate one block at a time, sample its coordinates uniformly on the unit circle, and fix $D_i^\top D_i=\diag(\alpha_{i,1},\alpha_{i,2})$, with $\alpha_{i,1}\ge\alpha_{i,2}>0$. As strengths vanish at fixed ratios, CE measures average squared leakage into the inactive block. Subject to these self-signals,
\begin{equation}
\boxed{\min_{D_1,D_2}\|D_1^\top D_2\|_F^2=\alpha_{1,2}\alpha_{2,2}.}
\label{eq:block-min}
\end{equation}
The planes must share a line. The cheapest arrangement shares their weak axes and separates their strong ones:
\begin{equation}
\begin{aligned}
D_1&=[\sqrt{\alpha_{1,1}}e_1\quad\sqrt{\alpha_{1,2}}e_3],\\
D_2&=[\sqrt{\alpha_{2,1}}e_2\quad\sqrt{\alpha_{2,2}}e_3].
\end{aligned}
\label{eq:block-solution}
\end{equation}
Both features remain two-dimensional: weak modes are shared, not deleted (Appendix~\ref{app:blocks}). When internal modes are also learned, retaining a mode gains predictive accuracy but may increase interference elsewhere. Selection within features and superposition between features therefore compete for the same dimensions.
\label{main:end}